We check the index set and the supplied class in that theorem
explicitly. Put $N_K=2Nd$ and $N'_K=cd_s$. These integers are
coprime by (K5), since $d\mid F$. Section~0.1 of \cite{kato1999}
requires every index $m$ with $N_K\mid m$ and $(m,N'_K)=1$.
The enlarged definition (K8) supplies all these classes. They
originate in the integral $j_*$ group over
$\mathbf Z[\zeta_m,1/2]$ and map to full cohomology over
$\mathbf Z[\zeta_m,1/N_K]$. The varying construction support
$S_{Dm}$ is not an instruction to invert all its primes in this
last scheme. Apply the fixed lattice map $2^{c_E}$ from (K4) if
needed, identically at every level.

For a prime $\ell\nmid N'_K$, if $\ell\mid m$ then
$S_{Dm\ell}=S_{Dm}$, so Proposition~8.12 gives norm factor one.
If $\ell\nmid m$, the condition $N_K\mid m$ implies
$\ell\nmid2Nd$, and hence
$S_{Dm\ell}\setminus S_{Dm}=\{\ell\}$. The equivariance identity
before (K9) gives the factor
$P_\ell(\ell^{-1}\sigma_\ell^{-1})=E_\ell(\chi)$, with
$P_\ell(X)=1-a_\ell\chi(\ell)X+\ell X^2$. These are exactly the
two cases of (0.1.1) in \cite{kato1999}. Smoothing primes are
excluded by $N'_K$ and need not be inverted.

For the class in Section~0.7 of that source take its parameter
$d=1$; this parameter is distinct from our odd conductor part $d$.
For every $n\geq1$,
\[
 S_{DN_K2^n}=S_{D2^n}=\operatorname{prime}(2Nd).
\]
Proposition~8.12 therefore gives
\[
 \operatorname{cor}_{\mathbf Q(\zeta_{DN_K2^n})/
                          \mathbf Q(\zeta_{D2^n})}
             z_{DN_K2^n}=z_{D2^n}.
\]
Transitivity of corestriction, through the two intermediate fields
in the common upper field, now yields
\[
 \operatorname{cor}_{\mathbf Q(\zeta_{N_K2^n})/
                          \mathbf Q(\zeta_{2^n})}y_{N_K2^n}
 =\operatorname{cor}_{\mathbf Q(\zeta_{D2^n})/
                          \mathbf Q(\zeta_{2^n})}
                         (z_{D2^n}\otimes e_\chi)
 =y_{2^n}.
\]
No Euler factor occurs because the supports are equal. No field
degree occurs because this is corestriction of the actual upper
class, not corestriction of a restriction. These equalities commute
with transition maps, and $n\geq1$ is a cofinal tail. Write
$\xi_K$ for the Section~0.7 class, to distinguish it from the
modular-symbol combination $\xi$. It is exactly the full-cohomology
image of $y_{2^\infty}$, multiplied by $2^{c_E}$ if the fixed
lattice map was needed. That factor is a unit at every height-one
prime under consideration, so its coordinate ideal is precisely
$J_{\mathfrak p}$. The preceding
nonvanishing argument supplies a nontorsion class on each sign
component, purity gives $w=-1$, and oddness gives condition
0.6(iv$_{\mathfrak p}$). At the height-one primes used here, $2$ is
invertible and the localized Iwasawa algebra is a DVR. Theorem 0.8
then bounds the coordinate ideal of this actual retained class. Its
strong residual and integral-cokernel assumptions concern primes
containing $2$, which are not used here.

We also account for the difference between the two published
cohomology conventions. Write $C_j$ for the $j_*$ Iwasawa complex of
\cite[Section 12.1]{kato}, and $C_S$ for full cochains on the open
scheme with the odd support inverted, as in \cite[Section 0.5]{kato1999}.
Work in one rational sign component and localize at $\mathfrak p$.
Their quotient is the direct sum of the odd local complexes (K16).
Each differential has nonzero determinant, so its degree-one
cohomology vanishes. Hence $H^1(C_j)\simeq H^1(C_S)$ and
$H^2(C_j)\hookrightarrow H^2(C_S)$. At an odd prime $q$, the
unramified local term has cohomological dimension one after inverting
$2$. The localization triangle identifies the degree-two cohomology
of (K16) with full local $H^2(\Q_q,T\chi)$ in Iwasawa cohomology.
The map from $H^2(C_S)$ is the usual local restriction. Exactness
therefore identifies $H^2(C_j)$ with the kernel of all those odd
local restrictions. Intersecting with the kernel of restriction at
$2$ gives
\[
 \ker\bigl(H^2(C_j)\to H^2_{\rm loc,2}\bigr)
 \simeq
 \ker\bigl(H^2(C_S)\to\!\bigoplus_{q\in S}H^2_{\rm loc,q}\bigr).
\]
This is precisely the comparison of the subscript-zero groups in
the two sources. The $H^1$ identification preserves the actual
Euler-system class and its coordinate ideal $J$. Thus Theorem 0.8
gives (K15) in the convention used here. The argument for (K16)
below is local and does not use (K15). We do not invoke the
odd-prime-only Perrin--Riou argument or an unchecked version of
Rubin's theorem at $2$.

